\section{El ge\'ofono como transductor: modelo y respuesta}
\label{sec:geofono}

Esta secci\'on trata el transductor en detalle porque de \'el se derivan casi todas las
especificaciones de la electr\'onica desarrollada.

\subsection{Modelo mecánico}

El ge\'ofono de bobina m\'ovil es un transductor \emph{inercial}: una masa suspendida por un
resorte, solidaria a una bobina que se mueve dentro del campo de un im\'an permanente
fijado a la carcasa. La carcasa sigue al suelo; la masa, por inercia, no. La tensi\'on de
salida se genera por la \textbf{velocidad relativa} entre bobina e im\'an.

Sea $x_g$ el desplazamiento del suelo (y por tanto de la carcasa) y $x_m$ el de la masa.
Definiendo el desplazamiento relativo $z=x_m-x_g$, la segunda ley de Newton sobre la masa
suspendida da

\begin{equation}
m\ddot{z}+c\dot{z}+kz=-m\ddot{x}_g
\quad\Longleftrightarrow\quad
\ddot{z}+2\zeta\omega_n\dot{z}+\omega_n^{2}z=-\ddot{x}_g,
\label{eq:geo-mecanico}
\end{equation}

con $\omega_n=\sqrt{k/m}$ la frecuencia natural y $\zeta=c/(2\sqrt{km})$ el factor de
amortiguamiento. La tensi\'on generada es proporcional a la velocidad relativa:

\begin{equation}
V_{\mathrm{out}}(t)=G_0\,\dot{z}(t),
\label{eq:geo-fem}
\end{equation}

donde $G_0$ es la sensibilidad, en \si{\volt\per\meter\per\second}.

\subsection{Las tres formas de la funci\'on de transferencia}
\label{sec:geo-tres-formas}

De \eqref{eq:geo-mecanico} y \eqref{eq:geo-fem} se obtiene la funci\'on de transferencia,
cuya \textbf{forma depende de qu\'e variable del movimiento del suelo se tome como
entrada}. Esta distinci\'on es importante y conviene fijarla expl\'icitamente, porque las
tres formas circulan en la literatura y difieren en un factor $s$:

\begin{align}
\text{entrada aceleraci\'on } \ddot{x}_g:\quad
H_a(s)&=\frac{V_{\mathrm{out}}(s)}{\ddot{X}_g(s)}
=\frac{-G_0\,s}{s^{2}+2\zeta\omega_n s+\omega_n^{2}},
\label{eq:geo-acel}\\[0.4em]
\text{entrada velocidad } \dot{x}_g:\quad
H_v(s)&=\frac{V_{\mathrm{out}}(s)}{\dot{X}_g(s)}
=\frac{-G_0\,s^{2}}{s^{2}+2\zeta\omega_n s+\omega_n^{2}},
\label{eq:geo-vel}\\[0.4em]
\text{entrada desplazamiento } x_g:\quad
H_d(s)&=\frac{V_{\mathrm{out}}(s)}{X_g(s)}
=\frac{-G_0\,s^{3}}{s^{2}+2\zeta\omega_n s+\omega_n^{2}}.
\label{eq:geo-desp}
\end{align}

\textbf{En este trabajo se adopta la forma en aceleraci\'on}, Ec.~\eqref{eq:geo-acel},
como descripci\'on de referencia del transductor, siguiendo el tratamiento de Ma
\emph{et al.}~\cite{Ma2023}, cuya topolog\'ia de acondicionamiento es la base de la cadena
anal\'ogica implementada. La raz\'on es pr\'actica: en la forma en aceleraci\'on el ge\'ofono se
comporta como un \textbf{filtro pasa-banda de segundo orden}, y el problema de dise\~no
---extender su banda \'util hacia frecuencias bajas--- se formula de manera directa como
un problema de modificaci\'on de los polos de esa funci\'on, tal como se desarrolla en la
Secci\'on~\ref{sec:topologias}.

Conviene observar que las tres formas describen el mismo dispositivo f\'isico: se
relacionan por $H_v=s\,H_a$ y $H_d=s\,H_v$, ya que en el dominio de Laplace derivar
equivale a multiplicar por $s$. La forma en velocidad, Ec.~\eqref{eq:geo-vel}, es la que
justifica la afirmaci\'on habitual de que ``el ge\'ofono mide velocidad'': para
$\omega\gg\omega_n$ se tiene $H_v\to -G_0$, es decir, respuesta plana y salida
proporcional a la velocidad de part\'icula del suelo.

\nota{\textbf{Confirmar}: en el documento fij\'e la forma en aceleraci\'on como referencia,
por coherencia con Ma \emph{et al.} y porque es la que usa el an\'alisis del compensador.
Si en el trabajo de MATLAB (\texttt{SM24\_model.slx}, \texttt{compensadorKai.m}) la
entrada del modelo est\'a definida como velocidad, hay que unificar el criterio o
declarar expl\'icitamente en cada lugar cu\'al se usa, para que las ganancias no queden
desfasadas en un factor $\omega$.}

\subsection{El problema de la respuesta en baja frecuencia}

Escribiendo \eqref{eq:geo-acel} en frecuencia, la magnitud de la respuesta en aceleraci\'on
tiene un m\'aximo en torno a $\omega_n$ y decae a ambos lados. Por debajo de $\omega_n$ la
ca\'ida es de \SI{20}{\decibel\per decade} en aceleraci\'on (equivalentemente
\SI{40}{\decibel\per decade} en la forma en velocidad). El resultado pr\'actico es que
la se\~nal de baja frecuencia se atenúa hasta hundirse en el piso de ruido del
amplificador.

Este es el conflicto central del proyecto, y conviene enunciarlo con precisi\'on:

\begin{itemize}[itemsep=0.15em]
    \item La informaci\'on de las \textbf{capas profundas} reside en las componentes de
    \textbf{baja frecuencia} de la curva de dispersi\'on, porque la profundidad de
    penetraci\'on es proporcional a la longitud de onda.
    \item El objetivo de \SI{50}{\meter} de profundidad exige, por lo tanto, recuperar
    contenido espectral por debajo de la frecuencia natural del transductor.
    \item El ge\'ofono SM-24 tiene $f_n=\SI{10}{\hertz}$, es decir, atenúa precisamente
    la banda de mayor inter\'es para el objetivo del proyecto.
\end{itemize}

La literatura del vault clasifica los ge\'ofonos seg\'un $f_n$ y su aplicaci\'on t\'ipica:

\begin{center}\footnotesize
\begin{tabular}{@{}cll@{}}
\toprule
$f_n$ (\si{\hertz}) & Aplicaci\'on t\'ipica & Uso en ondas superficiales \\
\midrule
4{,}5 & Exploraci\'on profunda, pasiva & Pasivo (microtremores, bajas frecuencias) \\
10   & Multiuso & Activo y pasivo \\
14   & MASW est\'andar & Activo (maza) \\
28   & Alta resoluci\'on superficial & Activo, capas muy someras \\
\bottomrule
\end{tabular}
\end{center}

Park \emph{et al.}~\cite{Park1999} emplearon ge\'ofonos de \SI{14}{\hertz} en el trabajo
fundacional del MASW, y para el rango dominado por la maza (\SIrange{5}{50}{\hertz}) la
respuesta es plana con variaciones menores al \SI{1}{\percent} en amplitud y
\SI{2}{\degree} en fase, sin necesidad de correcci\'on instrumental. Para trabajo pasivo o
profundidades mayores a \SI{30}{\meter} la recomendaci\'on habitual es
$f_n\leq\SI{4.5}{\hertz}$.

\subsubsection{Justificaci\'on de la elecci\'on del transductor}

La elecci\'on del SM-24 de \SI{10}{\hertz} responde a la \textbf{disponibilidad del
transductor en la Facultad}: es el ge\'ofono con el que se cuenta. No fue, por lo tanto,
una elecci\'on libre de dise\~no sino una restricci\'on de partida.

Esto es importante para interpretar correctamente el trabajo desarrollado. Para el
objetivo de \SI{50}{\meter} de profundidad la literatura recomendar\'ia
$f_n\leq\SI{4.5}{\hertz}$, de modo que el transductor disponible atenúa precisamente la
banda que interesa. En consecuencia, \textbf{el trabajo de extensi\'on de banda descrito en
la Secci\'on~\ref{sec:topologias} no es un refinamiento opcional sino la condici\'on que
hace viable el objetivo del proyecto con el hardware disponible}. Es, adem\'as, una
situaci\'on representativa del contexto que el proyecto busca atender: la de un
laboratorio que debe obtener resultados \'utiles con el instrumental que tiene, en lugar
de especificar el \'optimo y comprarlo.

\subsection{Par\'ametros del SM-24}

\begin{center}\footnotesize
\begin{tabular}{@{}lll@{}}
\toprule
Par\'ametro & Valor & Observaci\'on \\
\midrule
Frecuencia natural $f_n$ & \SI{10}{\hertz} nominal & Define el l\'imite inferior sin compensar \\
Sensibilidad $G_0$ & $\approx\SI{28.8}{\volt\per\meter\per\second}$ & Determina el nivel de se\~nal \\
Amortiguamiento $\zeta$ & $\approx0{,}7$ con resistencia de shunt & Valor de respuesta m\'as plana \\
\bottomrule
\end{tabular}
\end{center}

\nota{\textbf{Verificar contra la hoja de datos y contra la medici\'on}: resistencia de
bobina, inductancia, masa m\'ovil, amortiguamiento a circuito abierto y distorsi\'on. Estos
valores existen medidos en el trabajo de MATLAB (\texttt{SM-24/}) pero no los volqu\'e
ac\'a; conviene completar la tabla con \emph{valores medidos} adem\'as de los nominales,
porque es exactamente el tipo de dato que una mesa evaluadora pide.}

\subsection{Acoplamiento con el suelo}

El modelo anterior supone $\dot{x}_{\mathrm{ge\acute{o}fono}}=\dot{x}_{\mathrm{suelo}}$, lo
que exige un acoplamiento mec\'anico adecuado. Un acoplamiento deficiente ---superficie
dura, ge\'ofono inclinado, suelo blando no compactado--- introduce una
\textbf{resonancia espuria} en la respuesta y altera amplitud y fase del registro. El
efecto es cr\'itico si se pretende estimar el amortiguamiento del suelo $D_S$ a partir de
la curva de atenuaci\'on, y no despreciable para la propia curva de dispersi\'on.

\subsection{Alternativa evaluada: aceler\'ometros MEMS}

Se consider\'o el uso de aceler\'ometros MEMS en lugar de ge\'ofonos. Sus ventajas son
respuesta plana desde continua, salida digital, robustez y bajo costo; su desventaja es
un piso de ruido mayor en baja frecuencia frente a un ge\'ofono de calidad. Para MASW
activo en el rango \SIrange{5}{100}{\hertz} son pr\'acticamente equivalentes.

Se descart\'o el camino MEMS por dos razones: el ge\'ofono SM-24 ya estaba disponible, y
---m\'as importante--- el problema de acondicionamiento anal\'ogico de un transductor
pasivo de baja tensi\'on es precisamente donde el proyecto pod\'ia aportar contenido de
ingenier\'ia electr\'onica, que es el objeto de un PFC de Ingenier\'ia Electr\'onica. Un nodo
basado en MEMS habr\'ia sido, en lo esencial, un problema de integraci\'on digital.

\nota{\textbf{Confirmar} que esta es efectivamente la raz\'on del descarte, o corregirla.
No est\'a registrada en el repositorio; la reconstru\'i a partir del contexto.}
